\documentclass[11pt,a4paper,fontset=fandol]{ctexart}

\usepackage[a4paper,top=2.4cm,bottom=2.6cm,left=2.35cm,right=2.35cm,headheight=18pt,headsep=0.55cm,footskip=26pt]{geometry}
\usepackage{amsmath,amssymb}
\usepackage{fancyhdr}
\usepackage{lastpage}
\usepackage{enumitem}
\usepackage{titlesec}
\usepackage{xcolor}
\usepackage{hyperref}
\usepackage{bookmark}
\usepackage[most]{tcolorbox}
\usepackage{setspace}
\usepackage{array}

\hypersetup{
  colorlinks=true,
  linkcolor=blue!50!black,
  urlcolor=blue!50!black
}

\setstretch{1.16}
\setlength{\parindent}{2em}
\setlength{\parskip}{0.32em}
\setlist[enumerate]{itemsep=0.35em, topsep=0.35em, leftmargin=2.3em}
\setlist[itemize]{itemsep=0.25em, topsep=0.25em, leftmargin=2em}

\definecolor{mainblue}{RGB}{36,83,140}
\definecolor{lightblue}{RGB}{241,246,252}
\definecolor{lightgray}{RGB}{246,246,246}
\definecolor{rulegray}{RGB}{110,118,128}
\definecolor{softgreen}{RGB}{236,247,240}
\definecolor{softyellow}{RGB}{254,249,229}

\titleformat{\section}
  {\Large\bfseries\color{mainblue}}
  {\thesection.}
  {0.45em}
  {}
  [\vspace{0.25em}\color{rulegray}\titlerule]
\titlespacing*{\section}{0pt}{1.2em}{0.8em}
\titleformat{\subsection}{\large\bfseries\color{mainblue}}{\thesubsection}{0.5em}{}

\pagestyle{fancy}
\fancyhf{}
\fancyhead[L]{\small 错位排列专题目录页}
\fancyhead[C]{\small 组合专题资料索引}
\fancyhead[R]{\small 刘文博}
\fancyfoot[C]{\small 第 \thepage 页 / 共 \pageref{LastPage} 页}
\renewcommand{\headrulewidth}{0.55pt}
\renewcommand{\footrulewidth}{0.35pt}

\newtcolorbox{goalbox}[1]{
  breakable,
  colback=lightblue,
  colframe=mainblue,
  boxrule=0.7pt,
  arc=1mm,
  title=\textbf{#1},
  fonttitle=\bfseries
}

\newtcolorbox{infobox}[1]{
  breakable,
  colback=softgreen,
  colframe=green!40!black,
  boxrule=0.65pt,
  arc=1mm,
  title=\textbf{#1},
  fonttitle=\bfseries
}

\newtcolorbox{warnbox}[1]{
  breakable,
  colback=softyellow,
  colframe=orange!60!black,
  boxrule=0.65pt,
  arc=1mm,
  title=\textbf{#1},
  fonttitle=\bfseries
}

\begin{document}

\thispagestyle{empty}
\begin{titlepage}
\centering
\vspace*{0.9cm}

\begin{tcolorbox}[
  enhanced,
  width=0.92\textwidth,
  colback=white,
  colframe=mainblue,
  boxrule=1pt,
  arc=0mm,
  left=6mm,
  right=6mm,
  top=8mm,
  bottom=8mm
]
\centering

{\fontsize{24}{30}\selectfont\bfseries 错位排列专题目录页\par}
\vspace{0.65cm}
{\fontsize{17}{22}\selectfont 讲义、题单与周测的使用说明\par}

\vspace{1cm}
\begin{tcolorbox}[colback=lightblue,colframe=mainblue,width=0.84\textwidth,boxrule=1pt]
\centering
{\large 适合对象：已经掌握排列组合基础、准备系统提升“位置受限计数”的学生}\\[0.35em]
{\normalsize 本目录页帮助把单份资料串成一个完整的小专题训练包}
\end{tcolorbox}

\vspace{1.0cm}
\renewcommand{\arraystretch}{1.42}
\begin{tabular}{>{\bfseries}p{3.5cm}p{8cm}}
专题名称： & 错位排列 \\
资料组成： & 课堂讲义、提高训练题单、周测 A 卷、周测 B 卷 \\
使用方式： & 先学讲义，再做题单，最后用 A/B 周测检查掌握情况 \\
编译方式： & XeLaTeX \\
\end{tabular}

\vspace{0.9cm}
{\large \today\par}
\end{tcolorbox}
\end{titlepage}

\tableofcontents
\newpage

\section{专题包包含哪些资料}

\begin{goalbox}{四份核心资料}
\begin{enumerate}
  \item 课堂讲义：帮助建立概念、公式和常见模型；
  \item 提高训练题单：用于课后巩固与变式提升；
  \item 周测 A 卷：检查基础掌握是否稳定；
  \item 周测 B 卷：检查综合迁移与较难变式是否真正会做。
\end{enumerate}
\end{goalbox}

\begin{infobox}{对应文件名}
\begin{itemize}
  \item 讲义：\texttt{derangements\_handout.tex} / \texttt{derangements\_handout.pdf}
  \item 提高训练题单：\texttt{derangements\_advanced\_problem\_set.tex} / \texttt{derangements\_advanced\_problem\_set.pdf}
  \item 周测 A 卷：\texttt{derangements\_weekly\_test\_A.tex} / \texttt{derangements\_weekly\_test\_A.pdf}
  \item 周测 B 卷：\texttt{derangements\_weekly\_test\_B.tex} / \texttt{derangements\_weekly\_test\_B.pdf}
\end{itemize}
\end{infobox}

\section{建议学习顺序}

\begin{goalbox}{推荐顺序}
\begin{enumerate}
  \item 先完整学习讲义，重点吃透“什么是错位排列”“递推公式怎么来”“恰好若干个在原位怎么转化”；
  \item 再完成提高训练题单，先独立做，再对照详细解答订正；
  \item 接着做周测 A 卷，检查基础题型是否已经稳定；
  \item 最后做周测 B 卷，检查对容斥、递推和综合变式的掌握程度。
\end{enumerate}
\end{goalbox}

\begin{warnbox}{不建议的做法}
\begin{itemize}
  \item 只背公式，不理解为什么是 $D_n=(n-1)(D_{n-1}+D_{n-2})$；
  \item 一上来就做 B 卷，这样容易把“不会做”误以为“题太难”；
  \item 做完题只看答案数字，不对照解答过程总结错误原因。
\end{itemize}
\end{warnbox}

\section{每份资料主要解决什么问题}

\subsection{课堂讲义}
讲义适合第一次系统学习这个专题，主要解决以下问题：
\begin{itemize}
  \item 错位排列的定义到底是什么；
  \item 小规模情形怎么枚举；
  \item 递推公式从哪里来；
  \item 容斥公式和错位排列之间是什么关系；
  \item “至少一个在原位”“恰好几个在原位”怎么转化。
\end{itemize}

\subsection{提高训练题单}
题单适合第二阶段训练，主要用来把方法做熟。它更强调：
\begin{itemize}
  \item 识别题型；
  \item 选择公式还是容斥；
  \item 处理“部分对象受限”的变式；
  \item 避开高频陷阱，比如“恰好有 $n-1$ 个在原位”这种不可能情形。
\end{itemize}

\subsection{周测 A 卷}
A 卷偏基础，适合检查：
\begin{itemize}
  \item 标准错位排列是否会求；
  \item 恰好若干个在原位是否会转化为 $\binom{n}{k}D_{n-k}$；
  \item 至少若干个在原位是否会用补集法；
  \item 简单部分位置受限题是否会用容斥。
\end{itemize}

\subsection{周测 B 卷}
B 卷偏综合，适合检查：
\begin{itemize}
  \item 递推与容斥能否灵活切换；
  \item 复杂“前若干个对象受限”的题是否会写通项；
  \item 分类讨论时是否会遗漏不可能情况；
  \item 结果算出后能否用数量级检查答案是否合理。
\end{itemize}

\section{一个推荐的 4 次训练安排}

\begin{goalbox}{4 次完成一个小专题}
\begin{enumerate}
  \item 第 1 次：学习讲义前半部分，完成讲义中的例题与小练习；
  \item 第 2 次：学习讲义后半部分，完成提高训练题单的前半部分；
  \item 第 3 次：完成提高训练题单后半部分，并集中订正错题；
  \item 第 4 次：做周测 A 卷或 B 卷，作为阶段性检查。
\end{enumerate}
\end{goalbox}

\begin{infobox}{如果孩子基础较好，可以这样压缩}
\begin{itemize}
  \item 第 1 天：讲义 + 题单基础题；
  \item 第 2 天：题单变式题 + 错题订正；
  \item 第 3 天：周测 A 卷；
  \item 第 4 天：周测 B 卷。
\end{itemize}
\end{infobox}

\section{专题学习后的自检清单}

\begin{goalbox}{如果下面 8 条都能做到，说明这个专题基本过关}
\begin{enumerate}
  \item 我能准确说出什么是标准错位排列；
  \item 我知道 $D_1,D_2,D_3,D_4,D_5,D_6$ 的值；
  \item 我能解释递推公式为什么成立；
  \item 我会用容斥原理计算 $D_n$ 的小规模值；
  \item 我会处理“恰好有 $k$ 个在原位”；
  \item 我会处理“至少有 $k$ 个在原位”；
  \item 我会处理“只有前若干个对象受限”的变式；
  \item 我知道“恰好有 $n-1$ 个在原位”是不可能的。
\end{enumerate}
\end{goalbox}

\section{给家长的使用建议}

\begin{warnbox}{更有效的带练方式}
\begin{itemize}
  \item 不要一开始就给答案，先让孩子口头说“这题属于哪一类”；
  \item 如果卡住，可以只提示“想想补集法”或“想想是不是部分位置受限”；
  \item 订正时重点看思路是否清楚，而不只是最后数字是否正确；
  \item 错题建议隔两三天再重做一次，检验是否真正掌握。
\end{itemize}
\end{warnbox}

\begin{infobox}{和后续专题的衔接}
错位排列学完以后，最自然的下一步就是继续学习\textbf{容斥原理}。因为很多“部分对象不能在原位”的变式题，本质上已经在用容斥原理了。把这两个专题连起来学，效果会更好。
\end{infobox}

\end{document}
